\documentclass[a4paper,
fontsize=11pt,
%headings=small,
oneside,
numbers=noperiodatend,
parskip=half-,
bibliography=totoc,
final
]{scrartcl}

\usepackage[babel]{csquotes}
\usepackage{synttree}
\usepackage{graphicx}
\setkeys{Gin}{width=.4\textwidth} %default pics size

\graphicspath{{./plots/}}
\usepackage[ngerman]{babel}
\usepackage[T1]{fontenc}
%\usepackage{amsmath}
\usepackage[utf8x]{inputenc}
\usepackage [hyphens]{url}
\usepackage{booktabs} 
\usepackage[left=2.4cm,right=2.4cm,top=2.3cm,bottom=2cm,includeheadfoot]{geometry}
\usepackage[labelformat=empty]{caption} % option 'labelformat=empty]' to surpress adding "Abbildung 1:" or "Figure 1" before each caption / use parameter '\captionsetup{labelformat=empty}' instead to change this for just one caption
\usepackage{eurosym}
\usepackage{multirow}
\usepackage[ngerman]{varioref}
\setcapindent{1em}
\renewcommand{\labelitemi}{--}
\usepackage{paralist}
\usepackage{pdfpages}
\usepackage{lscape}
\usepackage{float}
\usepackage{acronym}
\usepackage{eurosym}
\usepackage{longtable,lscape}
\usepackage{mathpazo}
\usepackage[normalem]{ulem} %emphasize weiterhin kursiv
\usepackage[flushmargin,ragged]{footmisc} % left align footnote
\usepackage{ccicons} 
\setcapindent{0pt} % no indentation in captions
\usepackage{xurl} % Breaks URLs

%%%% fancy LIBREAS URL color 
\usepackage{xcolor}
\definecolor{libreas}{RGB}{112,0,0}

\usepackage{listings}

\urlstyle{same}  % don't use monospace font for urls

\usepackage[fleqn]{amsmath}

%adjust fontsize for part

\usepackage{sectsty}
\partfont{\large}

%Das BibTeX-Zeichen mit \BibTeX setzen:
\def\symbol#1{\char #1\relax}
\def\bsl{{\tt\symbol{'134}}}
\def\BibTeX{{\rm B\kern-.05em{\sc i\kern-.025em b}\kern-.08em
    T\kern-.1667em\lower.7ex\hbox{E}\kern-.125emX}}

\usepackage{fancyhdr}
\fancyhf{}
\pagestyle{fancyplain}
\fancyhead[R]{\thepage}

% make sure bookmarks are created eventough sections are not numbered!
% uncommend if sections are numbered (bookmarks created by default)
\makeatletter
\renewcommand\@seccntformat[1]{}
\makeatother

% typo setup
\clubpenalty = 10000
\widowpenalty = 10000
\displaywidowpenalty = 10000

\usepackage{hyperxmp}
\usepackage[colorlinks, linkcolor=black,citecolor=black, urlcolor=libreas,
breaklinks= true,bookmarks=true,bookmarksopen=true]{hyperref}
\usepackage{breakurl}

%meta
%meta

\fancyhead[L]{C. Lenz\\ %author
LIBREAS. Library Ideas, 48 (2026). % journal, issue, volume.
\href{https://doi.org/10.18452/37902}{\color{black}https://doi.org/10.18452/37902}
{}} % doi
\fancyhead[R]{\thepage} %page number
\fancyfoot[L] {\ccLogo \ccAttribution\ \href{https://creativecommons.org/licenses/by/4.0/}{\color{black}Creative Commons BY 4.0}}  %licence
\fancyfoot[R] {ISSN: 1860-7950}

\title{\LARGE{Expert*inneninterviews als Methode bibliothekswissenschaftlicher Forschung}}% title
\author{Charlotte Lenz} % author

\setcounter{page}{1}

\hypersetup{%
      pdftitle={Expert*inneninterviews als Methode bibliothekswissenschaftlicher Forschung},
      pdfauthor={Charlotte Lenz},
      pdfcopyright={CC BY 4.0 International},
      pdfsubject={LIBREAS. Library Ideas, 48 (2026)},
      pdfkeywords={Interviews, Methoden, Forschung, Bibliothekswissenschaft, Expert*inneninterviews, interviews, methodologies, research, library science},
      pdflicenseurl={https://creativecommons.org/licenses/by/4.0/},
      pdfcontacturl={http://libreas.eu},
      pdfurl={https://doi.org/10.18452/37902},
      pdfdoi={10.18452/37902},
      pdflang={de},
      pdfmetalang={de}
     }



\date{}
\begin{document}

\maketitle
\thispagestyle{fancyplain} 

%abstracts
\begin{abstract}
\noindent
\textbf{Kurzfassung}: Basis für das Schreiben von Hochschulschriften ist
nicht nur das Finden einer passenden Forschungsfrage, sondern vor allem
auch das Festlegen der Methodik, mit der die Daten, die für die
Beantwortung der Frage benötigt werden, erhoben werden sollen. Gerade im
Rahmen einer Dissertation wird der Entwicklung des methodischen
Vorgehens sogar ein ganzes Kapitel gewidmet -- dem sogenannten
Forschungsdesign. Die Anzahl der Instrumente, mit denen Forschungsdaten
ermittelt werden können, ist groß, wobei zwischen quantitativen und
qualitativen Methoden unterschieden werden kann. Die Methodik steht
dabei in Wechselwirkung mit der zu behandelnden Forschungsfrage.
Letztere legt fest, welche der Forschungsmethoden sich am besten für die
Beantwortung der Frage eignet. In vielen Fällen eignen sich geplante
Gespräche mit Fachleuten, die in dem Metier des
Untersuchungsgegenstandes tätig sind. Was man unter den sogenannten
Expert*inneninterviews versteht, wie sie durchgeführt und ausgewertet
werden und welchen Nutzen sie im Bereich der Bibliothekswissenschaft
haben, sollen die folgenden Ausführungen zeigen.

\begin{center}\rule{0.5\linewidth}{0.5pt}\end{center}

\noindent\textbf{Abstract}: The basis for an academic thesis does not consist
solely of an appropriate research question. The definition of the
methodology by which the required data are to be collected is necessary
as well. In the context of a dissertation, the development of the
methodological approach is even given an entire chapter---the so-called
research design. The number of instruments available for collecting
research data is huge, with a distinction being made between
quantitative and qualitative methods. Methodology is closely
interrelated with the research question under investigation, as the
latter determines which research methods are best suited to answering
the question. In many cases, planned interviews with experts who work in
the professional field related to the subject of investigation are
particularly suitable. The following discussion will explain what is
meant by so-called \enquote{expert interviews}, how they are conducted and
analyzed, and what benefits they offer in the field of library and
information science.
\end{abstract}

%body
\newpage

Die Autorin hat sich für das eigene Promotionsprojekt neben einer
systematischen Literaturanalyse auch für Expert*inneninterviews als
Forschungsmethode entschieden. Als Vorbereitung auf die damit
verbundenen Gespräche soll dieser Essay dienen, um sich mit der Methode
selbst besser vertraut zu machen. Der besseren Lesbarkeit willen,
erfolgen alle weiteren personenbezogenen Bezeichnungen in diesem Aufsatz
in der männlichen Form, welche dabei aber keinerlei wertende Bedeutung
hat und alle Geschlechter gleichermaßen repräsentiert.

\section{1. Definition}\label{definition}

Ein Experteninterview ist eine erweiterte Form des sogenannten
Leitfaden-Interviews, also eine qualitative Befragung anhand eines vorab
erstellten Fragenkatalogs. Die für das Interview herangezogenen
Interviewpartner fungieren dabei als Quelle, aus der Wissenschaftler
wertvolle Erkenntnisse über den von ihnen zu untersuchenden Gegenstand
herausarbeiten können. Voraussetzung dabei ist, dass der Befragte eine
fachliche Kompetenz bezüglich der Forschungsthematik besitzt.
(Vergleiche Werner 2013, S. 142 ff.) Er sollte also über ein
überdurchschnittliches Fachwissen im Hinblick auf das Thema verfügen.
Seine Befragung dient als methodisches Instrument, dieses besondere
Wissen an die Oberfläche zu heben. (Vergleiche Gläser und Laudel 2010,
S. 11) Ziel der Befragung ist es, die Bedeutung des Standpunktes und der
Perspektive der Experten für den Kontext der Untersuchung
herauszuarbeiten und auf Grundlage ihrer Aussagen spezifische Annahmen
herzuleiten, um die zentrale Forschungsfrage zu beantworten. Warum sich
besonders solche Personen für ein Interview eignen, verdeutlicht Lars
Rinsdorf in seinem Beitrag \emph{Qualitative Befragungen}. (Vergleiche
Rinsdorf 2013, S. 67 ff.) Seiner Auffassung nach sind diese Befragten,
die er als \emph{Beforschte} bezeichnet, tatsächlich \emph{Experten}, da
sie im Arbeitsalltag in dem zu untersuchenden Fachgebiet tätig sind. Als
sogenannte \emph{relevante Akteure} können sie eine komplexe und
umfassende Bewertung zum Thema abgeben, die sich neben der angewandten
Praxis auch aus einer entsprechenden Vorbildung und persönlichen
Erfahrungen speist.

\section{2. Arbeitsschritte}\label{arbeitsschritte}

Ein Experteninterview sollte in fünf Schritten erfolgen:

\subsection{1. Die Vorüberlegungen}\label{die-voruxfcberlegungen}

Zunächst gilt es, die Wer-, Wie- und Wo-Fragen zu beantworten. Es wird
also festgelegt, welche Personengruppe, unter Anwendung welcher
Hilfsmittel und an welchen Orten beziehungsweise über welche Kanäle
befragt werden soll. Im Bereich der Bibliothekswissenschaft können dies
sowohl Hochschulangehörige, wie Lehrende und Studierende, als auch
Praktiker, vom Bibliotheksleiter bis zum Fachangestellten, sein, wenn
sich die Untersuchung im Rahmen von einzelnen Arbeitsbereichen in einer
Bibliothek dreht. Aber auch Bibliotheksbenutzende stellen eine, wenn
nicht sogar die primären, Experten für bibliothekswissenschaftliche
Forschung dar, wenn zum Beispiel eine Bewertung von bibliothekarischen
Dienstleistungen im Mittelpunkt der Forschung steht. Ausschlaggebend
ist, wie bereits erwähnt, immer die Forschungsfrage, die die zu
befragende Fokusgruppe meist in sich impliziert.

Der Ort der Befragung bestimmt die zu verwendenden Hilfsmittel. Findet
das Interview in Persona statt, reicht womöglich ein Leitfaden, etwas
zum Schreiben und/oder ein Diktiergerät. Findet das Gespräch online
statt, sollte sichergestellt werden, dass sowohl der Interviewführende
als auch der Befragte jeweils über die technischen Voraussetzungen (PC,
Software für das Durchführen von Videokonferenzen) verfügen. Generell
muss ein Termin gefunden werden, der es beiden Parteien erlaubt, das
Interview ohne Zeitdruck und andere Ablenkungen zu absolvieren. Nur dann
wird es am Ende für beide Seiten, aber im Besonderen für den
Forschenden, erfolgreich verlaufen.

Hat man sich auf die Fokusgruppe sowie die Art der Befragung festgelegt,
gilt es, passende Interviewpartner ausfindig zu machen und sie um einen
Interviewtermin zu bitten.

\subsection{2. Die Vorbereitung des
Interviews}\label{die-vorbereitung-des-interviews}

Im Anschluss sollte der Ablauf des Interviews geplant werden. Dies
erfolgt in der Regel anhand eines Leitfadens mit zielorientierten
Fragen, die als roter Faden durch das Gespräch führen sollen. In einer
qualitativen Befragung dient er dazu, das zu untersuchende Themenfeld zu
strukturieren und in der Befragungssituation selbst als
Orientierungshilfe, um nicht vom Thema abzuschweifen. Wenn dieser im
Vorfeld geteilt wird, ermöglicht er beiden Seiten festzustellen, ob alle
relevanten Punkte, die für die Beantwortung der Forschungsfrage wichtig
sind, besprochen wurden. Grundsätzlich sollten die Fragen so formuliert
werden, dass sich das Gegenüber zum Antworten ermuntert fühlt. Auch ist
es sinnvoll, weiterführende Fragen zu stellen sowie dem Interviewpartner
die Freiheit zu geben, eigene Punkte anzusprechen. Die Fragen selbst
werden dem zu Interviewenden im Vorfeld zur Verfügung gestellt, damit
dieser sich auf das Gespräch und gegebenenfalls auf weitere, für ihn
wichtige Themen vorbereiten kann. Die Möglichkeit und selbstverständlich
auch Zeit für einen offenen Dialog unabhängig von den Leitfragen sollten
auch angedacht werden. Für den Aufbau dieser Art von Leitfäden empfehlen
Bogner et al.~die Einteilung in mehrere Einzelabschnitte, welche jeweils
maximal drei Leitfragen enthalten sollten, die in jedem Fall abgefragt
werden sollten. Sollte im Gespräch bemerkt werden, dass diese Fragen
nicht ausführlich genug beantwortet wurden, können zusätzlich
vertiefende Fragen gestellt werden. (Vergleiche Bogner u.~a. 2014, S. 27
ff.)

\subsection{3. Die Durchführung des
Interviews}\label{die-durchfuxfchrung-des-interviews}

Nun gilt es, das Interview durchzuführen. Damit dies in Gänze gelingt
und das damit verbundene Ziel, nämlich Informationen, die zur
Beantwortung der Forschungsfrage notwendig sind, zu erhalten, erreicht
wird, muss der Ablauf gut geplant werden. Idealerweise kommt es während
der Unterhaltung mit der Zeit zu einem für beide Parteien angenehmen
Fachgespräch, welches zwei Phasen umfassen sollte. Zu Beginn, in der als
\emph{Eröffnungssequenz} bezeichneten Phase, sollte durch den
Forschenden eine Situation geschaffen werden, in der sich der Befragte
wohl, aber vor allem in seiner Expertise ernst genommen fühlt. Die
Vermittlung der eigenen Qualifikation bezüglich des Forschungsthemas,
aber auch ein sympathisches Auftreten können zu einer positiven
Entwicklung dieses Stadiums des Interviews beitragen. (Vergleiche Kaiser
2014, S. 93 ff.) Um dies zu erreichen, könnte vorab das
Forschungsprojekt vorgestellt, das Ziel des Gesprächs verdeutlicht,
sowie Gründe aufzeigt werden, warum der Interviewführende gerade im
jeweiligen Gesprächspartner einen Sachkundigen bezüglich des
Forschungsthemas sieht. Darüber hinaus wäre es sinnvoll, über den Ablauf
der Befragung zu informieren, eine Einwilligung zur eventuellen Aufnahme
des Gesprächs einzuholen und die weitere Verwendung der durch das
Gespräch entstehenden Informationen zu erklären. Ferner sollte über die
Art und Weise belehrt werden, wie die im Zuge des Interviews erhobenen
Forschungsdaten gespeichert werden sollen. Nicht zuletzt sollte die
Option unterbreitet werden, die Befragung vertraulich zu halten.
(Vergleiche Kaiser 2014, S. 92). Zum anderen könnte auch der
Interviewpartner durch gezielte Fragen zur eigenen Person und zum
Bildungs- und Berufsweg einbezogen werden. Solche eher unverfänglichen
Fragen können in der Regel leicht beantwortet werden und schaffen beiden
Parteien Sicherheit für die zweite Phase des Gesprächs. Diese zweite
Phase wird mehr Zeit erfordern als das eben erläuterte Warm-up, denn
hier wird all jenes abgefragt, was zur Beantwortung der Forschungsfrage
relevant ist. Für beide Akteure und für das Herausarbeiten der besten
Informationen bezüglich der Forschungsfrage wäre es optimal, wenn es an
dieser Stelle zu einem Gespräch auf Augenhöhe, in einem, von Robert
Kaiser als \emph{Modus} \emph{der Argumentation und Diskussion}
bezeichnet, kommen würde. (Vergleiche Kaiser 2014, S. 94) Den Abschluss
der Durchführung des Interviews bildet eine Dokumentation über dessen
Ablauf. Da sich neben den geäußerten Inhalten die Interviewsituation
selbst auf das Ergebnis der erhobenen Forschungsdaten nachhaltig
auswirkt, sollte diese neben den Hauptaussagen des Gesprächs
festgehalten werden, um bei der Auswertung ebenfalls berücksichtigt
werden zu können. Gläser et al.~empfiehlt daher, Informationen zur
Motivation des Interviews, dessen Rahmenbedingungen und den Ablauf des
Interviews zu protokollieren. (Vergleiche Gläser und Laudel 2010, S. 192
ff.)

Das bestmögliche Gelingen von Experteninterviews erreicht man, wenn man
vor der Durchführung der eigentlichen Interviews über einen sogenannten
Pre-Test die Tauglichkeit des Fragenkatalogs sowie die
Interviewsituation im Beisein eines Beispielbefragten erprobt. Solche
Tests sind fester Bestandteil der Forschungsmethode und sollten
ernsthaft durchgeführt und ebenfalls dokumentiert werden, denn an dieser
Stelle wird sich zeigen, ob Fragen gegebenenfalls umformuliert oder der
Ablauf des Gesprächs selbst verändert werden müssen. (Vergleiche Kaiser
2014, S. 82 ff.)

\subsection{4. Die Transkription und Auswertung des
Interviews}\label{die-transkription-und-auswertung-des-interviews}

Für die Dokumentation des Interviews ist eine vollständige Transkription
des gesamten Gesprächs unerlässlich, damit auch zu einem späteren
Zeitpunkt nachvollzogen werden kann, wie der Forschende auf die für die
Beantwortung der Forschungsfrage benötigten Ergebnisse gekommen ist. Bei
der Verschriftlichung sollten nur die reinen Aussagen übernommen werden.
Füllwörter dürfen ignoriert werden. Unvollständige Sätze können ohne
Sinnveränderung geschlossen werden. Sollte es dem Verständnis dienen
beziehungsweise einen besonderen Einfluss auf die Forschungsfrage haben,
können spontane, aber auch bewusste Pausen entsprechend gekennzeichnet
werden. Nach der Übertragung der Interviews werden diese dem
Interviewpartner zur Überprüfung zur Verfügung gestellt. Die Auswertung
der Interviews erfolgt dann auf Grundlage der Transkription und des nach
dem Interview erstellten Gesprächsprotokolls, wobei versucht wird, die
Kernaussagen herauszuarbeiten. Die Gewinnung der für die Forschungsfrage
relevanten Informationen erfolgt letztlich durch eine vergleichende
Auswertung aller wesentlichen Daten, die über Interviews mit
verschiedenen Experten, die zur gleichen Thematik befragt wurden,
zustande gekommen sind. Dies zeigt deutlich, dass für eine objektive
Erhebung von Forschungsdaten und einer ebensolchen Beantwortung der
Forschungsfrage niemals nur ein Experteninterview erfolgen darf. Es gibt
natürlich keine Standardvorgabe dafür, wie viele Experten befragt werden
sollten, aber die Zahl sollte so gewählt sein, dass die Aussagen für das
Thema aussagekräftig sind. Die wissenschaftliche Analyse der Daten
erfolgt über das Zusammentragen aller wichtigen inhaltlichen
Gesichtspunkte, welche über ein Kategoriensystem ausgewertet werden. Die
Kategorien werden im Anschluss gewichtet und die darin enthaltenen
Informationen grob als Pro- oder Kontraargument bezüglich der
Beantwortung der Forschungsfrage eingeordnet, falls dies die
Fragestellung zulässt. So können inhaltlich zusammenhängende Fakten, die
aber von unterschiedlichen Interviewpartnern geäußert wurden,
zusammengefasst werden. Tatsächlich können über dieses Verfahren auch
Fehlinterpretationen oder -einschätzungen der Interviewpartner
aufgedeckt und ausgewiesen werden. (Vergleiche Bogner u.~a. 2014, S. 72
ff.)

\subsection{5. Die Verschriftlichung der erhobenen Daten im Hinblick auf
den
Untersuchungsgegenstand}\label{die-verschriftlichung-der-erhobenen-daten-im-hinblick-auf-den-untersuchungsgegenstand}

Anhand dieser aufbereiteten Daten kann das Ergebnis der Befragung mit
Bezug auf die Forschungsfrage sinnvoll verschriftlicht werden. Für die
Wahrung der Neutralität der Untersuchung ist es unabdingbar, bei diesem
Arbeitsschritt auch auf persönliche Grenzen beim Interview, auf Probleme
bei einem oder mehreren der Arbeitsschritte, auf die Forschung stark
beeinflussende Vorkommnisse bei der Befragung sowie die zur Durchführung
und Auswertung verwendete technische Hilfsmittel und die damit in
Verbindung eventuell auftretenden Hürden einzugehen. Dies wird die
Erhebung der Daten vollständig dokumentieren und anderen Forschenden,
die sich für eigene Projekte der gleichen Forschungsmethode bedienen
möchten, eine enorme Hilfe sein.

\section{3. Der Nutzen von Experteninterviews als
Forschungsmethode in der
Bibliothekswissenschaft}\label{der-nutzen-von-experteninterviews-als-forschungsmethode-in-der-bibliothekswissenschaft}

Experteninterviews dienen dazu, das Forschungsthema in Gänze zu
beleuchten sowie aus dem Wissen und dem Erfahrungsschatz der Befragten
ein vollständiges Bild bezüglich der zentralen Forschungsfrage zu
zeichnen. Damit lassen sich Meinungen, Einschätzungen, aber auch
Einwände oder Bedenken ermitteln, auf die der Forschende von selbst
möglicherweise gar nicht gekommen wäre. Eigene Erfahrungen, wenn auch
\textquotesingle subjektiv\textquotesingle{} wahrgenommen, sind darüber
hinaus ein Instrument, theoretische Annahmen auf ihre Richtigkeit und
Übertragbarkeit auf die Praxis hin zu überprüfen. Hierin besteht der
wesentliche Nutzen dieser Methode für die bibliothekswissenschaftliche
Forschung. Bogner et al. fassen dies wie folgt perfekt zusammen:
\enquote{Das Besondere am Expertenwissen besteht nicht nur in dessen
besonderer Reflexivität, Kohärenz oder Gewissheit, sondern auch
insbesondere darin, dass dieses Wissen in besonderer Weise praxiswirksam
und damit orientierungs- und handlungsleitend für andere Akteure wird.}
(Bogner u.~a. 2014, S. 13 ff.) Dieser Einfluss umfasst dabei nicht nur
das eigene Arbeitsumfeld der Experten, sondern kann sich auch auf
angrenzende oder sogar fachübergreifende Bereiche erstrecken, was den
Stellenwert von Experteninterviews für die verschiedensten
Forschungsprojekte noch erhöht. Für manche wissenschaftliche
Forschungsvorhaben sind sie sogar unverzichtbar. (Vergleiche Bogner
u.~a. 2014, S. 13 ff.)

\section{Literatur}\label{literatur}

Bogner, Alexander, Beate Littig und Wolfgang Menz. 2014.
\emph{Interviews mit Experten - Eine praxisorientierte Einführung}.
Qualitative Sozialforschung Lehrbuch. Springer VS.
\url{https://doi.org/10.1007/978-3-531-19416-5}

Gläser, Jochen, und Grit Laudel. 2010. \emph{Experteninterviews und
qualitative Inhaltsanalyse}. 4. Auflage. Lehrbuch. VS Verlag für
Sozialwissenschaften.

Kaiser, Robert. 2021. \emph{Qualitative Experteninterviews -
Konzeptionelle Grundlagen und praktische \linebreak Durchführung}. 2.,
aktualisierte Auflage. Elemente der Politik. Springer VS.
\url{https://doi.org/10.1007/978-3-658-30255-9}

Rinsdorf, Lars. 2013. \enquote{Qualitative Methoden}. In: \emph{Handbuch
Methoden der Bibliotheks- und Informationswissenschaft}, herausgegeben
von Konrad Umlauf, Simone Fühles-Ubach und Michael Seadle. De Gruyter
Saur. \url{https://doi.org/10.1515/9783110255546.64}

Werner, Petra. 2013. \enquote{Qualitative Befragungen}. In
\emph{Handbuch Methoden der Bibliotheks- und Informationswissenschaft},
herausgegeben von Konrad Umlauf, Simone Fühles-Ubach und Michael Seadle.
De Gruyter Saur. \url{https://doi.org/10.1515/9783110255546.128}

%autor
\begin{center}\rule{0.5\linewidth}{0.5pt}\end{center}

\textbf{Charlotte Lenz}, hat Bibliotheksmanagement an der
Fachhochschule Potsdam sowie Bibliotheks- und Informationswissenschaften
an der HU Berlin studiert. Nach einer langjährigen Anstellung als
wissenschaftliche Bibliothekarin in der Bibliothek des Deutschen
Historischen Museums betreut sie nun die Bibliothek der Gedenkstätte
Sachsenhausen.

\end{document}